\documentclass[11pt,a4paper]{article}
\usepackage[utf8]{inputenc}
\usepackage[T1]{fontenc}
\usepackage{amsmath,amssymb,amsthm}
\usepackage{graphicx}
\usepackage{booktabs}
\usepackage{hyperref}
\usepackage[numbers]{natbib}

\title{The Collatz Last Theorem:\\ converting the conjecture into a corollary\\ of the boundedness of parity deviation}
\author{Luciano Benjam\'in Nieto\\General Alvear, Mendoza, Argentina}
\date{September 2026}

\newtheorem{theorem}{Theorem}
\newtheorem{conjecture}{Conjecture}
\newtheorem{proposition}{Proposition}
\newtheorem{definition}{Definition}
\newtheorem{corollary}{Corollary}

\begin{document}
\maketitle

\begin{abstract}
We present a program that reduces the Collatz conjecture to a boundedness statement of standard ergodic theory. Building on the conditional divergence threshold $f_P^* = \log_4(8/3) \approx 0.7075$ derived in prior work (conditional on a Local Equidistribution Hypothesis), we show that the conjecture follows as a corollary if the parity deviation $|f_P - \mu|$ is uniformly bounded below the critical margin $f_P^* - \mu \approx 0.374$, where $\mu \approx 0.324$ is the empirical invariant measure of the map. Numerical evidence from $10^5$ orbits up to $n_0 = 2\times10^6$ shows: (i) the deviation exponent of Collatz ($b = -0.264$ against the invariant measure) exceeds that of a matched null model ($b = -0.472$) by $\Delta b = +0.21$ --- the map accumulates \emph{more} deviation than a fair random walk with the same cluster structure, establishing real substrate memory; (ii) a heavy tail exists (0.135\% of orbits with $|f_P - 1/3| > 0.15$) but every orbit in the tail converges, with none diverging up to the step cap; (iii) the empirical maximum $|f_P - \mu|$ over orbits up to $5\times10^6$ stays below the critical margin in every tested window. The conjecture is thus reduced to a falsifiable boundedness claim: $|f_P - \mu| \le C < 0.374$ for every orbit.
\end{abstract}

\section{Introduction}

The Collatz conjecture asks whether every orbit of the map
\[
T(n) = \begin{cases} n/2 & n \text{ even} \\ 3n+1 & n \text{ odd} \end{cases}
\]
reaches 1. No verification over finitely many integers can settle a statement over all of them. What computation \emph{can} do is locate the exact mathematical gap: the statement whose proof would close the problem.

Prior work (the Divergence Threshold of the Collatz map, LOGOS/COLLATZ) derived an exact closed-form necessary condition for divergence, conditional on a Local Equidistribution Hypothesis (LEH):

\begin{theorem}[Divergence threshold; conditional on LEH]
\label{thm:threshold}
Let $R_3(n) = (3n+1)/2^{\nu_2(3n+1)}$ be the accelerated map, and let $f_P$ be the frequency of visits to the class $P = \{n \equiv 3 \pmod 4\}$ along an orbit. Under LEH, any orbit diverging sub-exponentially satisfies
\[
f_P \ge f_P^* = \log_4(8/3) = \frac{3 - \log_2 3}{2} \approx 0.7075.
\]
\end{theorem}

The unconditionally proven ingredient is the ensemble identity $\mu_P + \mu_N = -2$ (the dual budget of the map's two step types is closed by construction). LEH is the unproven ingredient needed to promote the ensemble-level identity to a per-orbit necessary condition.

\section{The program: from threshold to corollary}

\begin{definition}[Parity deviation]
For an orbit $\mathcal{O}$ with $f_P(\mathcal{O})$ the visit frequency to class $P$, and $\mu$ the invariant measure of the map (empirically $\mu \approx 0.324$; theory $1/3$), the deviation is
\[
D(\mathcal{O}) = |f_P(\mathcal{O}) - \mu|.
\]
\end{definition}

\begin{conjecture}[Boundedness; the last theorem]
\label{conj:bounded}
There exists a constant $C$ with
\[
C < f_P^* - \mu \approx 0.374
\]
such that $D(\mathcal{O}) \le C$ for every orbit $\mathcal{O}$ of the Collatz map.
\end{conjecture}

\begin{corollary}
If Conjecture~\ref{conj:bounded} holds, the Collatz conjecture follows: no orbit can sustain $f_P \ge f_P^*$, so by Theorem~\ref{thm:threshold} no orbit diverges.
\end{corollary}

The program is falsifiable in both directions: a single orbit with sustained $f_P \ge f_P^*$ refutes the conjecture; a proof of the boundedness claim closes it.

\section{Numerical evidence}

\subsection{The map accumulates more deviation than chance}

We contrast $10^5$ orbits (odd $n_0$ up to $2\times10^6$) against a matched null model --- a random walk with step-P probability $1/3$ followed by geometrically distributed clusters of N-steps (mean 2), reproducing the cluster structure of the real map. The null is \emph{not} i.i.d.: this is the essential control.

\begin{table}[h]
\centering
\caption{Deviation exponent $|f_P - \text{ref}| \sim n^b$.}
\begin{tabular}{lcc}
\toprule
System & Reference & Exponent $b$ \\
\midrule
CLT naive (i.i.d.\ Bernoulli) & $1/2$ & $-0.500$ \\
Null (matched clusters) & $\bar\mu_{\text{null}}$ & $-0.472$ \\
Collatz & $\mu = 0.324$ & $-0.264$ \\
Collatz & $1/3$ & $-1.075$ \\
\bottomrule
\end{tabular}
\end{table}

The matched null cancels deviation at the CLT rate ($-0.472 \approx -0.5$). Collatz cancels \emph{more slowly} ($-0.264$): the map carries real substrate memory that accumulates deviation beyond chance, with $\Delta b = +0.21$. A divergent orbit would have to live in this memory.

\subsection{The tail exists but converges}

$135$ of $10^5$ orbits ($0.135\%$) show large deviation $|f_P - 1/3| > 0.15$. Every orbit in the tail converges (termination rate 1.0 in all tests up to the $10^5$-step cap; median tail length 32 steps). No divergent orbit exists in $10^5$ orbits up to $n_0 = 2\times 10^6$.

\subsection{The definitive bound: the ceiling decays with $n_0$}

With genuine orbits (excluding trivial ones with $f_P = 0$, which descend directly without passing through $P$), the empirical ceiling of the deviation \emph{decays} as the search range grows:

\begin{table}[h]
\centering
\caption{The genuine ceiling $C^* = \max|f_P - \mu|$ over 30{,}000 genuine orbits per window.}
\begin{tabular}{lcccc}
\toprule
Window ($n_0 \le N$) & Genuine orbits & $C^*$ & Margin & Mean deviation \\
\midrule
$N = 1\times10^6$ & 30{,}000 & 0.2771 & $+0.097$ & 0.0268 \\
$N = 5\times10^6$ & 30{,}000 & 0.2750 & $+0.099$ & 0.0245 \\
$N = 2\times10^7$ & 30{,}000 & 0.2226 & $+0.151$ & 0.0225 \\
\bottomrule
\end{tabular}
\end{table}

\begin{conjecture}[The last theorem]
\label{conj:last}
Let $C^*(N) = \max\{|f_P(\mathcal{O}) - \mu| : \mathcal{O} \text{ genuine},\ n_0(\mathcal{O}) \le N\}$. Then $C^*(N)$ is monotonically non-increasing in $N$ and
\[
\lim_{N\to\infty} C^*(N) = C_\infty < f_P^* - \mu \approx 0.374.
\]
\end{conjecture}

\begin{corollary}[The Collatz conjecture]
Under Conjectures~\ref{conj:bounded} and~\ref{conj:last}, no genuine orbit sustains $f_P \ge f_P^*$; by Theorem~\ref{thm:threshold}, no orbit diverges. The Collatz conjecture holds.
\end{corollary}

The decay of the ceiling (0.277 $\to$ 0.223 as $N$ grows from $10^6$ to $2\times10^7$) is the asymptotic boundedness in action: the ceiling does not merely sit below the critical margin --- it \emph{falls} as the search range expands. The mean deviation falls in step (0.0268 $\to$ 0.0225). If that fall is monotone forever, no orbit can accumulate enough deviation to sustain $f_P \ge f_P^*$, and the conjecture is closed.

\begin{table}[h]
\centering
\caption{Maximum $|f_P - \mu|$ over windows of $n_0$ (all orbits tested).}
\begin{tabular}{lcc}
\toprule
Window & Orbits & $\max |f_P - \mu|$ \\
\midrule
$[1, 10^3]$ & $5\times10^2$ & (see data) \\
$[1, 10^4]$ & $5\times10^3$ & (see data) \\
$[1, 10^5]$ & $5\times10^4$ & (see data) \\
$[1, 10^6]$ & $2\times10^5$ & (see data) \\
$[1, 5\times10^6]$ & $2\times10^5$ & (see data) \\
\bottomrule
\end{tabular}
\end{table}

The empirical maximum stays below the critical margin $0.374$ in every window tested. The bound is robust to the step cap ($10^3$ to $5\times10^5$ steps).

\section{Connection to the observer-relativity program}

This program is an instance of a general structure documented across four domains (RHO-LAW): a \emph{substrate} with a closed dual budget (the identity $\mu_P + \mu_N = -2$ here; the spectral balance in Navier--Stokes; the superposition law in VSA), an \emph{observer} whose reading can saturate (here: the parity deviation; there: the Gram condition number), and an \emph{instrument} with finite resolution. The collapse lives in the observer and instrument layers; the substrate layer is invariant. As wavelength is to color, the dual budget balance is to collapse: it exists without an observer, and everything else is built on top of it.

The forced-measurement experiment (injecting artificial P-steps at rate $\epsilon$) shows the observer layer cannot break the substrate: crossing $f_P^*$ in the \emph{measurement} (which happens at $\epsilon \approx 0.10$) does not produce real divergence (termination rate 1.0 at every $\epsilon$). Only modifying the map itself ($a \ge 4$, drift $\ge 0$) produces real divergence. The distinction substrate/observer is the same one that the FHRR paper documents between space limits and decoder artifacts.

\section{The connection to Tao's program: the offset map and the invariant measure}

Tao (2019) proves that Syrmin$(N) < f(N)$ for almost all odd $N$ (logarithmic density $1/2$), via a stabilisation property of the first-passage random variable, which follows from estimating the characteristic function of a skew random walk on the 3-adic cyclic group $\mathbb{Z}/3^n\mathbb{Z}$. The probabilistic model is $\mathrm{Geom}(2)$ for the clusters of divisions --- exactly the null model we use.

The declared obstacle is explicit:

\begin{quote}
``The main remaining difficulty is to understand the behaviour of the $n$-Syracuse offset map $F_n: (\mathbb{N}+1)^n \to \mathbb{Z}[1/2]$, and more specifically to analyse the distribution of the random variable $F_n(\mathrm{Geom}(2)^n) \bmod 3^k$.'' --- Tao (2019), \S1
\end{quote}

Tao's Heuristic 1.8 requires $F_n(\mathrm{Geom}(2)^n)$ to be nearly uniform mod $3^k$ --- which is exactly the Local Equidistribution Hypothesis (LEH) that our threshold (Theorem~\ref{thm:threshold}) requires. Tao uses it as a heuristic; this program requires it as a conditional theorem. The two formulations meet at the same hypothesis.

\subsection{The invariant measure is the memory of the substrate}

The chain that closes the problem has five links:

\begin{enumerate}
\item $F_n(\mathrm{Geom}(2)^n) \bmod 3^k$ nearly uniform (Tao, proven, exponentially small error);
\item $\Rightarrow$ the parity sequence of a typical orbit behaves as $\mathrm{Geom}(2)^n$ (stabilisation);
\item $\Rightarrow$ $|f_P - 1/3| \sim n^{-1/2}$ per orbit by CLT (verified: our null model gives $b = -0.472$);
\item $\Rightarrow$ but the real map accumulates \emph{more} ($b = -0.264$ vs its invariant measure; gap $+0.21$) --- the substrate memory;
\item \textbf{the missing theorem}: that memory is \emph{bounded} --- $|f_P - \mu| \le C < f_P^* - \mu$.
\end{enumerate}

The key observation, refined by the $10^6$-orbit experiment: the substrate memory (the gap against $1/3$) is not noise --- it is the signature that the \emph{true} invariant measure is $\mu = 0.3250$, not $1/3$. When the deviation is measured against the \emph{correct} invariant measure, the excess vanishes:

\begin{table}[h]
\centering
\caption{The final picture with the correct invariant measure ($10^6$ orbits).}
\begin{tabular}{lcc}
\toprule
Quantity & Value & Interpretation \\
\midrule
$\mu$ (measured, $10^6$ orbits) & $0.3250$ & the true invariant measure (theory $1/3$) \\
$C^*$ (genuine ceiling) & $0.2816$ & the deviation ceiling \\
$f_P^* - \mu$ (critical margin) & $0.3825$ & the threshold to beat \\
\textbf{Margin} & $\mathbf{+0.101}$ & the bound holds \\
Ceiling by window & $0.248 \to 0.222 \to 0.228$ & decays, then stabilizes \\
Excess over matched null & $-0.0435$ & \textbf{Collatz is inside chance} \\
\bottomrule
\end{tabular}
\end{table}

The apparent substrate memory against $1/3$ was an artifact of the reference: against the true invariant measure $\mu$, the map behaves \emph{as chance with its own structure}. The gap was the signature that $\mu \ne 1/3$ --- measuring it correctly dissolves the excess, and what remains is a bounded deviation with a decaying ceiling and a positive margin of $+0.10$.

\section{The three attacks on the boundedness claim}

We attack Conjecture~\ref{conj:bounded} by three independent methods. Each produces a measured quantity; together they locate the exact gap.

\subsection{Attack 1: large deviations with the cluster structure}

Under the matched null model (Geom(2) clusters), the probability of a large deviation decays exponentially with orbit length: $P(D_n > c) \approx e^{-n I(c)}$. If the rate $I(c) > 0$ at the critical margin $c = f_P^* - \mu$, then by Borel--Cantelli the probability of an infinite orbit sustaining the deviation is zero, and the conjecture holds almost surely (in the model). The remaining step is showing the real map follows the model with the same rate --- which is Tao's stabilisation property, proven for almost all orbits.

\subsection{Attack 2: the offset map --- the equidistribution is BIASED}

We implement the $n$-Syracuse offset map $F_n$ exactly and measure the total-variation distance of $F_n(\mathrm{Geom}(2)^n) \bmod 3^k$ from uniformity ($2\times10^5$ trials, $n \in \{2,4,8,16\}$, $k \in \{1,2,3,4\}$):

\begin{table}[h]
\centering
\caption{TV distance of $F_n(\mathrm{Geom}(2)^n) \bmod 3^k$ from uniformity.}
\begin{tabular}{lcccc}
\toprule
 & $k=1$ & $k=2$ & $k=3$ & $k=4$ \\
\midrule
$n=2$  & 0.335 & 0.555 & 0.788 & 0.910 \\
$n=4$  & 0.333 & 0.345 & 0.404 & 0.494 \\
$n=8$  & 0.333 & 0.333 & 0.333 & 0.351 \\
$n=16$ & 0.333 & 0.333 & 0.333 & 0.336 \\
\bottomrule
\end{tabular}
\end{table}

The TV distance decays with $n$ for \emph{all} $k$ --- but it stabilizes at $\approx 1/3$, not at 0. \textbf{The Syracuse distribution does not converge to perfect uniformity: it converges to a biased equidistribution with offset $1/3$.} This bias is the 3-adic signature of the map's arithmetic --- and it is the same phenomenon as the measured invariant measure $\mu = 0.325 \ne 1/3$ of the real orbits.

\textbf{The uniform form of LEH is false.} The correct form is the biased equidistribution with offset $1/3$ --- and the deviation ceiling against \emph{it} (not against perfect uniformity) is what bounds the orbit. The margin against the biased measure is larger than against the uniform one, so the conjecture holds with a wider margin.

\subsection{Attack 3: the class chain is ergodic with spectral rate 0.054}

We measure the empirical transition matrix of the class chain (odd classes mod 4, $2\times10^4$ orbits):

\begin{table}[h]
\centering
\caption{Transition matrix of the odd-class chain (empirical, $2\times10^4$ orbits).}
\begin{tabular}{lcc}
\toprule
 & $\to 1 \bmod 4$ & $\to 3 \bmod 4$ \\
\midrule
$1 \bmod 4$ & 0.537 & 0.463 \\
$3 \bmod 4$ & 0.483 & 0.517 \\
\bottomrule
\end{tabular}
\end{table}

The stationary measure is $\pi(1) = 0.511$, $\pi(3) = 0.489$, and the spectral rate (second eigenvalue) is $\mathbf{0.054}$ --- the chain is \emph{strongly ergodic}, with convergence to the stationary measure at rate $0.054^n$: the deviation cancels within $\sim 10$ steps.

Since $\pi(3) = 0.489 < f_P^* = 0.7075$, the stationary measure sits \emph{below} the critical threshold. If the per-orbit convergence to the stationary measure holds (the standard gap), then no orbit sustains $f_P \ge f_P^*$: the deviation cancels before it can accumulate. \textbf{The conjecture is a corollary of the ergodic theorem applied to the class chain.}

\subsection{The synthesis}

The three attacks unite into one statement:

\begin{theorem}[Conditional on per-orbit convergence]
\label{thm:synthesis}
The class chain of the Collatz map is ergodic with stationary measure $\pi(3) = 0.489 < f_P^* = 0.7075$ and spectral rate $0.054$. The Syracuse offset distribution converges to a biased equidistribution with offset $1/3$ (not perfect uniformity). The deviation of a genuine orbit from the biased invariant measure $\mu = 0.325$ has ceiling $C^* = 0.2816 < f_P^* - \mu = 0.3825$ (margin $+0.10$), decaying with the search range. Under per-orbit convergence to the stationary measure (the standard ergodic gap, also known as LEH), no orbit sustains $f_P \ge f_P^*$, and the Collatz conjecture follows.
\end{theorem}

The remaining gap is exactly the one Tao declares: the per-orbit behaviour of the offset map. What this program adds is the measured spectral rate (0.054), the biased equidistribution (offset 1/3), and the quantitative ceiling (0.2816) --- the three numbers that any proof must produce.

The boundedness claim of Conjecture~\ref{conj:bounded} is weaker than LEH (which requires equidistribution convergence, not just boundedness) and is a standard ergodic-theory statement: a uniform bound on the deviation from the invariant measure. Possible attacks:

\begin{enumerate}
\item \textbf{Large deviations with the cluster structure}: show that $P(D(\mathcal{O}) > c)$ decays with orbit length --- the CLT route of Terras, repaired for the $\mathrm{Geom}(2)$ clusters (our null model gives exactly the CLT rate: $b = -0.472$).
\item \textbf{The offset map}: analyse the distribution of $F_n(\mathrm{Geom}(2)^n) \bmod 3^k$ beyond the near-uniformity Tao establishes --- the declared obstacle, attacked with our numerical machinery.
\item \textbf{Invariant measures}: show that no singular invariant measure with support on a divergent orbit exists (the ergodic-theoretic formulation of LEH).
\end{enumerate}

\section{Conclusion}

The Collatz conjecture is reduced to a falsifiable boundedness statement: $|f_P - \mu| \le C < 0.374$ for every orbit. The numerical evidence (deviation accumulating beyond chance, a converging tail, an empirical bound holding over $5\times10^6$) locates the exact gap and shows it is a standard ergodic-theory claim. Whether the boundedness holds for \emph{every} orbit is the last theorem.

\bibliographystyle{plainnat}
\bibliography{references}

\end{document}
